\documentclass[11pt]{article}
\usepackage[margin=0.9in]{geometry}
\usepackage{amsmath,amssymb,mathtools}
\usepackage{xcolor}
\usepackage[hidelinks]{hyperref}
\usepackage{microtype}

\newcommand{\R}{\mathbb R}
\newcommand{\Sc}{\mathcal S}

\title{The Intermediate Schatten Gelfand-Number Gap\\
\large An explicit literature answer to arXiv:2011.06554}
\author{Automated mathematics run \texttt{fa\_banach\_001}, lane 14}
\date{22 July 2026}

\begin{document}
\maketitle

\begin{center}
\fcolorbox{black}{gray!8}{\parbox{0.88\textwidth}{
\textbf{Status: literature\_already\_answered.}
This note records an exact later-literature answer, not a new result.
}}
\end{center}

\section*{Original question}

Hinrichs, Prochno, and Vyb\'iral study the Gelfand numbers of the identity
\[
 c_n(\Sc_p^N\hookrightarrow\Sc_q^N),
 \qquad 1\leq p\leq2\leq q\leq\infty.
\]
In Section 4.6, p.~21 of the arXiv PDF, they state that their upper and lower
bounds match for small and very large codimension, and write:

\begin{quote}
``We leave it an open problem to find a good lower bound also in the
intermediate range
$c_{p,q}N^2\leq n\leq N^2-N_q+1$,'' where
$N_q=cN^{1+2/q}$.
\end{quote}

This is the exact signal extracted from arXiv:2011.06554 \cite{HPV}.

\section*{Explicit answering paper}

Prochno and Strzelecki \cite{PS} explicitly identify the connection. On p.~5
they say that their first Gelfand-number result ``closes a gap'' in the 2020
paper and that the intermediate-codimension upper bounds there are sharp.
Their Theorem A (p.~6) and Proposition 3.1 (p.~19) prove, for
$0<p\leq2\leq q\leq\infty$,
\[
 c_n(\Sc_p^N\hookrightarrow\Sc_q^N)
 \gtrsim N^{-1/2-1/p}\sqrt{N^2-n+1}
\]
whenever
\[
 (1-c)N^2\leq n\leq N^2-cN^{1+2/q}+1.
\]
This matches the upper bound in the original paper. For
$1\leq p\leq2$, their alternative proof gives lower-bound constant $1$.

Proposition 3.2 (p.~20) treats the complementary range:
\[
 c_n(\Sc_p^N\hookrightarrow\Sc_q^N)
 \gtrsim
 \min\left\{1,\frac{N^{3/2-1/p}}{\sqrt n}\right\},
 \qquad 1\leq n\leq(1-c)N^2.
\]
Together Propositions 3.1 and 3.2 cover the full intermediate strip in the
source question and give the matching lower asymptotics sought there.

\section*{Identification and scope}

This is an exact \texttt{literature\_already\_answered} match: the supporting
authors cite the original paper, call out its gap, and prove the matching
lower bounds. No portion is promoted as an original proof. The identification
is limited to the source's $1\leq p\leq2\leq q\leq\infty$ intermediate-range
question; separate conjectures for $2\leq p\leq q$ and quasi-Banach regimes in
the supporting paper are outside this packet.

\begin{thebibliography}{9}
\bibitem{HPV}
A. Hinrichs, J. Prochno, and J. Vyb\'iral,
\emph{Gelfand numbers of embeddings of Schatten classes},
arXiv:\href{https://arxiv.org/abs/2011.06554}{2011.06554} (2020),
Section 4.6, p.~21.

\bibitem{PS}
J. Prochno and M. Strzelecki,
\emph{Approximation, Gelfand, and Kolmogorov numbers of Schatten class
embeddings},
arXiv:\href{https://arxiv.org/abs/2103.13050}{2103.13050} (2021),
Theorem A and Propositions 3.1--3.2.
\end{thebibliography}

\end{document}
